\section{Reference values}\label{sec:references}

A published price accurate enough to test a high-order solver does not exist for this problem, so we
build our own. The three references below are independent of the collocation solver; how independent
they are of each other is discussed at the end of the section.

\subsection{Exact price of the geometric twin}\label{sec:geoexact}

Write $W^{(1)} = \rho W^{(2)} + \sqrt{1-\rho^2}\,W^{\perp}$ and let $Y_T = \frac1T\int_0^T\ln S_u\,\dd u$,
so that $G_T=e^{Y_T}$ and $Y_T=\zeta_T$. From the variance equation,
$\sqrt{v}\,\dd W^{(2)} = (\dd v - \kappa(\theta-v)\,\dd t)/\xi$, and integration by parts with $a(0)=1$,
$a(T)=0$ gives $\int_0^T a\,\dd v = -v_0 + \frac1T\int_0^T v\,\dd t$. The stochastic integral against
$W^{(2)}$ can therefore be written in terms of the variance path alone, and conditionally on that path
$Y_T$ is Gaussian. Its characteristic function is
\begin{equation}\label{eq:geochf}
\E\bigl[e^{iuY_T}\bigr] = \exp\bigl(iu\Lambda + \mathcal A(T) + \mathcal B(T)v_0\bigr),\qquad
\Lambda = \ln S_0 + \frac{rT}{2} - \frac\rho\xi\Bigl(v_0 + \frac{\kappa\theta T}{2}\Bigr),
\end{equation}
where, in time-to-maturity $\tau$ and with $a=\tau/T$,
\begin{equation}\label{eq:georiccati}
\frac{\dd\mathcal B}{\dd\tau} = \tfrac12\xi^2\mathcal B^2 - \kappa\mathcal B + c(\tau),\qquad
\frac{\dd\mathcal A}{\dd\tau} = \kappa\theta\mathcal B,\qquad \mathcal A(0)=\mathcal B(0)=0,
\end{equation}
\begin{equation*}
c(\tau) = iu\Bigl[-\frac a2 + \frac\rho\xi\Bigl(\frac1T + \kappa a\Bigr)\Bigr] - \frac{u^2}{2}(1-\rho^2)a^2 .
\end{equation*}
The Riccati equation has a time-dependent coefficient and we integrate it numerically (an eighth-order
Runge--Kutta method with relative tolerance $10^{-12}$), for all quadrature nodes at once. The price then
follows from the damped Fourier inversion of \citet{CarrMadan1999}, with damping $\alpha=1.5$ and
512 Gauss--Legendre nodes on $[0,200]$. The same representation underlies the semi-closed form of
\citet{KimWee2014}.

We checked this computation in three ways. In the limit $\xi\to0$, $\rho=0$, $v_0=\theta$ it reproduces the
Black--Scholes formula for the geometric Asian call with a relative error below $1.1\times10^{-11}$ for all
combinations of $r\in\{0,0.04,0.1\}$, $T\in\{0.25,1,5\}$ and dividend yield $q\in\{0,0.02\}$. For the full
model the price changes by less than $5\times10^{-11}$ when the damping is varied between $0.5$ and $3$ or
the quadrature is doubled, at maturities up to five years. And a direct simulation of the geometric payoff agrees with it within one standard error at every
maturity we tried ($T=0.25$, $3$ and $5$, two million paths each). For the base parameters
the prices are $13.3650818582$, $6.7095975217$ and $2.7174933218$ for $K=90$, $100$ and $110$.

\subsection{The lower bound of Fusai and Kyriakou under Heston}\label{sec:fk}

For any $\lambda\in\mathbb R$ the payoff satisfies $(A_T/T-K)^+ \ge (A_T/T-K)\,\ind\{Y_T>\lambda\}$, hence
\begin{equation}\label{eq:lb}
C(K) \ge \LB(\lambda) := e^{-rT}\,\E\bigl[(A_T/T-K)\,\ind\{Y_T>\lambda\}\bigr]
\end{equation}
for every $\lambda$. The idea of conditioning the arithmetic average on the geometric one goes back to
\citet{Curran1994} and \citet{RogersShi1995}. \citet{FusaiKyriakou2016} maximise \eqref{eq:lb} over
$\lambda$, which gives the maximal lower bound $\MLB = \LB(\lambda^\ast)$ with
$\E[A_T/T \mid Y_T=\lambda^\ast]=K$, and compute it by Fourier inversion in $\lambda$. With a damping
$\delta>0$,
\begin{equation}\label{eq:Phi}
\Phi(u;\delta) = \int_{\mathbb R} e^{(iu+\delta)\lambda}\,\LB(\lambda)\,\dd\lambda
= \frac{e^{-rT}}{iu+\delta}\Bigl[\frac1T\int_0^T\varphi_t(u-i\delta)\,\dd t - K\,\varphi_0(u-i\delta)\Bigr],
\end{equation}
\begin{equation}\label{eq:phis}
\varphi_t(\upsilon) = \E\bigl[S_t\,e^{i\upsilon Y_T}\bigr],\qquad \varphi_0(\upsilon) = \E\bigl[e^{i\upsilon Y_T}\bigr],
\end{equation}
and $\LB(\lambda) = \frac{e^{-\delta\lambda}}{\pi}\int_0^\infty \mathrm{Re}\bigl[e^{-iu\lambda}\Phi(u;\delta)\bigr]\dd u$.
(In the notation of \citet{FusaiKyriakou2016}, $\varphi_t=\phi_{t,t,T}(-i/2,\cdot)$ and
$\varphi_0=\phi_{T,T,T}(0,\cdot)$; see their equation (46).) For L\'evy models they give $\varphi_t$ in
closed form. For affine stochastic volatility models with a continuous average they describe the route,
via the joint characteristic function of variance, log-price and its integral and iterated conditioning,
but do not give formulas.

For the Heston model we compute $\varphi_t$ as follows; the details are in
Appendix~\ref{app:transform}. Let $X=\ln S$ and $c=i\upsilon/T$, so that
$S_te^{i\upsilon Y_T} = \exp\bigl(X_t + c\int_0^TX_u\,\dd u\bigr)$. Conditioning first on
$\mathcal F_t$, the expectation of $\exp(c\int_t^TX_u\,\dd u)$ is exponential-affine in $(X_t,v_t)$. The
product with $e^{X_t}$ and $\exp(c\int_0^tX_u\,\dd u)$ is again exponential-affine, and a second
conditioning on $\mathcal F_0$ finishes the computation. In time-to-maturity $\tau$ both stages solve the
same system
\begin{equation}\label{eq:twostage}
\frac{\dd B_x}{\dd\tau} = c,\qquad
\frac{\dd B_v}{\dd\tau} = \tfrac12B_x^2 - \tfrac12B_x + \rho\xi B_xB_v + \tfrac12\xi^2B_v^2 - \kappa B_v,\qquad
\frac{\dd \mathcal A}{\dd\tau} = rB_x + \kappa\theta B_v ,
\end{equation}
starting from zero, except that $B_x$ jumps by one at $\tau=T-t$, where the factor $e^{X_t}$ enters.
Hence $B_x(\tau) = c\tau + \ind\{\tau\ge T-t\}$, and
\begin{equation*}
\varphi_t(\upsilon) = \exp\bigl(\mathcal A(T) + (1+i\upsilon)\ln S_0 + B_v(T)\,v_0\bigr).
\end{equation*}
The first stage, $\tau\in[0,T-t]$, is the same for every $t$ and is integrated once with output at the
required points; the second stage is integrated separately for each quadrature node in $t$. Both are
vectorised over the Fourier variable and integrated numerically, so no branch of a complex logarithm has
to be tracked. We use 48 Gauss--Legendre nodes in $t$, 1024 nodes in $u$ on $[0,U]$ with
$U = 40/\sqrt{\bar vT/3}$ and $\bar v=\max(v_0,\theta)$, and $\delta=1$. The maximum over $\lambda$ is found by
a bounded scalar search, which is cheap because $\Phi$ is computed once and each evaluation of
$\LB(\lambda)$ is a single weighted sum. One parameter set with three strikes takes a few seconds in
Python.

\begin{table}[t]
\centering
\caption{Reproduction of Table~10 of \citet{FusaiKyriakou2016}: maximal lower bound (MLB) for the
Gaussian model, obtained with the Heston machinery in the limit $\xi\to0$ ($\xi=10^{-9}$, $\rho=0$,
$v_0=\theta=\sigma^2$). The last column is Linetsky's spectral price from the same table, for scale.}
\label{tab:fk10}
\footnotesize\setlength{\tabcolsep}{4.5pt}
\begin{tabular}{ccccc cccc}
\toprule
$S_0$ & $K$ & $r$ & $\sigma$ & $T$ & MLB (FK) & MLB (ours) & difference & Linetsky price\\
\midrule
2   & 2 & 0.02   & 0.10 & 1 & 0.055985 & 0.05598508 & $+7.9\times10^{-8}$ & 0.0559860415\\
2   & 2 & 0.18   & 0.30 & 1 & 0.218366 & 0.21836625 & $+2.5\times10^{-7}$ & 0.2183875466\\
2   & 2 & 0.0125 & 0.25 & 2 & 0.172226 & 0.17222617 & $+1.7\times10^{-7}$ & 0.1722687410\\
1.9 & 2 & 0.05   & 0.50 & 1 & 0.193060 & 0.19305997 & $-2.6\times10^{-8}$ & 0.1931737903\\
2   & 2 & 0.05   & 0.50 & 1 & 0.246298 & 0.24629819 & $+1.9\times10^{-7}$ & 0.2464156905\\
2.1 & 2 & 0.05   & 0.50 & 1 & 0.306094 & 0.30609360 & $-4.0\times10^{-7}$ & 0.3062203648\\
2   & 2 & 0.05   & 0.50 & 2 & 0.349779 & 0.34977879 & $-2.1\times10^{-7}$ & 0.3500952190\\
\bottomrule
\end{tabular}
\end{table}

We validated the implementation in four independent ways before using it.
\begin{itemize}
\item Without the jump in $B_x$ and at $t=T$, the system \eqref{eq:twostage} computes $\varphi_0$, the
characteristic function of $Y_T$. The geometric reference of Section~\ref{sec:geoexact} computes the same
quantity by a different derivation (it eliminates the log-price through the variance path), and the two
agree to $1.3\times10^{-15}$ on five parameter sets, including $T=5$ and a nonzero dividend yield.
\item For $\upsilon=0$ the system must return $\E[S_t]=S_0e^{rt}$ exactly, which it does to machine
precision.
\item In the Gaussian limit, $\LB(\lambda)$ has a closed form, because $(X_t,Y_T)$ is jointly Gaussian.
Our Fourier values agree with it to about $10^{-13}$ for all $\lambda$ tested, and the maxima reproduce all
seven values of Table~10 of \citet{FusaiKyriakou2016} to the six printed digits
(Table~\ref{tab:fk10}).
\item For the Heston parameter sets of their Table~9, a direct Monte Carlo estimate of $\LB(\lambda^\ast)$
at our $\lambda^\ast$ agrees with our Fourier value within two standard errors in nine runs
(three step sizes per set, standard errors $1.4$--$5.4\times10^{-4}$).
\end{itemize}
The bound is also stable: changing $\delta$ between $0.25$ and $2$ or doubling either quadrature moved it
by at most $1.4\times10^{-10}$ in every parameter set where we tried it, far below the Monte Carlo
uncertainty. Fusai and Kyriakou also give an upper bound for the error of the MLB. We did not
compute it: it requires conditional second moments of the average, and the values they report for
discretely monitored Heston options, $0.013$ to $0.049$, are far wider than the accuracy we need.

\subsection{Monte Carlo with the bound as control variate}\label{sec:e8}

Since the bound is exact for the continuous average, it can be used to correct a simulation, an idea
already used by \citet[Section~5]{FusaiKyriakou2016}. Along each path the difference
\begin{equation}\label{eq:gap}
\begin{aligned}
D &= e^{-rT}\Bigl[(A_T/T-K)^+ - (A_T/T-K)\,\ind\{Y_T>\lambda^\ast\}\Bigr]\\
  &= e^{-rT}\Bigl[(A_T/T-K)^+\ind\{Y_T\le\lambda^\ast\} + (K-A_T/T)^+\ind\{Y_T>\lambda^\ast\}\Bigr]
\end{aligned}
\end{equation}
is nonnegative and small: it is the part of the payoff that the bound misses, and it is nonzero only on
paths where the geometric average points the wrong way. We estimate
\begin{equation*}
C(K) = \MLB + \E[D]
\end{equation*}
by simulating $\E[D]$ with the quadratic-exponential scheme of \citet{Andersen2008}, antithetic normals,
and the trapezoidal rule for $A_T$ and $Y_T$. The time-discretisation bias of this estimator is the bias
of $\E[D]$ alone, and since $D$ is small, so is its bias.

For the base parameters $\E[D]$ is $0.0043$, $0.0036$ and $0.0038$ for $K=90,100,110$. We ran seven
independent batches with 32, 64 (eight million paths each) and 128, 256, 512, 1024 and 2048 (four million
paths each) time steps. The batches are mutually consistent ($\chi^2$ of 1.2, 9.4 and 7.2 on six degrees
of freedom), and a weighted regression of the estimates on the step size gives slopes within $1.8$
standard errors of zero. We therefore pool them. The result, which we call the reference price
throughout, is shown in Table~\ref{tab:ref}; we quote it with an uncertainty of $\pm3\times10^{-5}$,
which also covers the difference between the pooled value and the intercept of the regression. On the same
paths, a control variate based on the exact geometric price \citep{KemnaVorst1990} leaves a standard
error eight to ten times larger. The geometric control variate is the usual choice for arithmetic Asian
options under stochastic volatility, used recently by \citet{CustersKarbachFriesen2026} for discrete
monitoring, so the gain from the bound may be of use beyond the continuous case, although we have not
tested that. An estimate of that kind from eight million paths with 512 steps, which
is how we computed the reference before the bound was available, is listed for comparison; it differs
from the new reference by $1.3$ to $1.6$ of its own standard errors.

\begin{table}[t]
\centering
\caption{Reference prices for the base parameters ($\kappa=1$, $\theta=0.09$, $\xi=1$, $\rho=-0.3$,
$v_0=0.09$, $S_0=100$, $r=0.04$, $T=1$). MC: pooled estimate $\MLB+\E[D]$ from 36 million QE paths.}
\label{tab:ref}
\small
\begin{tabular}{lccc}
\toprule
& $K=90$ & $K=100$ & $K=110$\\
\midrule
Maximal lower bound (Section~\ref{sec:fk}) & 13.755959 & 7.028864 & 2.963683\\
$\E[D]$ & 0.004338 & 0.003634 & 0.003811\\
Reference price $\MLB+\E[D]$ & 13.760297 & 7.032498 & 2.967494\\
\quad standard error & $1.7\times10^{-5}$ & $1.5\times10^{-5}$ & $1.7\times10^{-5}$\\
\quad intercept of regression in $\Delta t$ & 13.760288 & 7.032471 & 2.967472\\
Geometric control variate only, 8M paths & 13.759848 & 7.032017 & 2.967033\\
\quad standard error & $3.5\times10^{-4}$ & $3.2\times10^{-4}$ & $2.8\times10^{-4}$\\
\bottomrule
\end{tabular}
\end{table}

At long maturities $\E[D]$ is larger ($0.014$--$0.035$ for $T=3$ and $5$), and there its bias does become
visible: at $T=3$ the estimate with 256 steps differs from those with 1024 and 2048 steps by up to
$6\times10^{-4}$. For these maturities we use 1024 steps or more.

\paragraph{How independent are the references?} The Monte Carlo reference and the bound do not use the
collocation solver in any way, and the bound's main ingredient, $\varphi_t$, is validated against three
separate checks. They do share with the geometric reference the affine structure of the model, numerical
integration of Riccati equations and Fourier inversion, and the reference price in Table~\ref{tab:ref}
depends on the bound directly. What the solver is compared with is therefore independent of the solver, but
not of the affine machinery. In Section~\ref{sec:published} we add a simulation with a different
time-stepping scheme and no control variate at all, which shares nothing with our other computations except
the model.
